\subsection{Topological D-brane generators}
For either threefold $X$, let $L_a$ denote its ordered primitive divisor
basis: $H_1,\ldots,H_6$ on the resolution and $K_1,K_2,K_3$ on the
smooth side. Put $c_a=c_2(X)\cdot L_a$ and
$\kappa_{abc}=L_aL_bL_c$. Work in untwisted topological $K^0(X)$,
with Dirac pairing minus the Euler pairing. Choose generators
\begin{equation}
 \mathsf m_0=[\mathcal O_X],\qquad
 \mathsf m_a=[\mathcal O_X(L_a)]-[\mathcal O_X],\qquad
 \mathsf e_0=-[\mathcal O_{\rm pt}],\qquad
 \operatorname{ch}(\mathsf e_a)=(0,0,\beta_a,0),
 \label{eq:x9-d6-k-generators}
\end{equation}
where $L_a\cdot\beta_b=\delta_{ab}$. Sans-serif symbols denote charge
generators, distinct from the exceptional curves $e_1,e_2$.
The sign of $\mathsf e_0$ gives
normalized period $+1$, consistently with the earlier central-charge
convention. In particular $\mathsf m_0$ has Mukai charge
$(1,0,c_2(X)/24,0)$, not the unshifted cohomological unit.

These generators form the full integral charge lattice. On the smooth
side, the vertex matrices of the seven-ray fan polytope and its polar
both have nonzero Smith entries $(1,1,1,1)$. The same property was
established for the resolution. The fundamental-group and Brauer-group
formulas therefore give torsion-free integral cohomology on both sides
\cite[Corollaries~1.9 and~3.9]{Batyrev:2005Integral}.
The untwisted Atiyah--Hirzebruch spectral sequence degenerates: its
differentials are rationally zero and its groups are torsion-free.
Its free associated graded groups are $H^0,H^2,H^4,H^6$
\cite[Section~3.1]{Maldacena:2001KCharges}.
Thus every primitive $H^4$ class has a rank-zero, first-Chern-zero
$K$-theory lift. Its index is $\int_X\operatorname{ch}_3\in\ZZ$,
so a point class removes its third Chern character. This constructs
the $\mathsf e_a$ in \eqref{eq:x9-d6-k-generators}; the remaining generators
lift the primitive $H^0,H^2,H^6$ bases. No BPS representative of every
generator is assumed. In particular the smooth-side products of two
divisors generate an index-ten sublattice of $H^4$, not all of $H^4$.
